%------------------------------------------------------------------------------%

\usepackage{apresentacao}
\usepackage{tabto}

\makeatletter
\graphicspath  {{./figs/}}
\def\input@path{{./figs/}}
\makeatother

%------------------------------------------------------------------------------%
\newcommand{\mytitle}   {Teste da Comparação}
\newcommand{\mysubtitle}{Séries Numéricas -- Testes de Convergência}
\title   {\mytitle}
\subtitle{\mysubtitle}
\author  {\large Luis Alberto D'Afonseca}
\date    {}

%------------------------------------------------------------------------------%
\begin{document}

%------------------------------------------------------------------------------%
\begin{frame}
    \titlepage
\end{frame}

%------------------------------------------------------------------------------%
\section{Teste da Comparação}

%------------------------------------------------------------------------------%
\definecolor{c1}{rgb}{0,.7,0}
\definecolor{c2}{rgb}{0,0,.7}
\definecolor{c3}{rgb}{.7,0,0}
\begin{frame}{Teste da Comparação}
  Sejam três séries com termos não negativos
  \[
    {\color{c1} \sum_{n=1}^{\infty} p_n} \qquad
    {\color{c2} \sum_{n=1}^{\infty} a_n} \qquad
    {\color{c3} \sum_{n=1}^{\infty} q_n}
  \]
  com \( \quad {\color{c1}p_n} \;\leq\; {\color{c2}a_n} \;\leq\; {\color{c3}q_n} \quad \) para todo $n>N$

  \pause
  \hspace{\baselineskip}
  \begin{enumerate}
    \item se \tabto{8mm} \(\color{c3} \sum q_n \) convergir \tabto{40mm} então  \(\color{c2} \sum a_n \) converge \\[6mm]\pause
    \item se \tabto{8mm} \(\color{c1} \sum p_n \) divergir  \tabto{40mm} então  \(\color{c2} \sum a_n \) diverge
  \end{enumerate}
\end{frame}

%------------------------------------------------------------------------------%
\begin{frame}{Exemplo}
  Para analisar a série do \emph{inverso do fatorial}
  \[
    \sum_{n=1}^{\infty} \frac{1}{\,n!\,}
    = 1             \;+\;
      \frac{1}{2}   \;+\;
      \frac{1}{3!}  \;+\;
      \frac{1}{4!}  \;+\;
      \frac{1}{5!}  \;+\;
      \cdots
  \]
  \pause
  comparamos com a série geométrica
  \[ \sum_{n=1}^{\infty} \frac{1}{2^{n-1}}
    = 1              \;+\;
      \frac{1}{2}    \;+\;
      \frac{1}{2^2}  \;+\;
      \frac{1}{2^3}  \;+\;
      \frac{1}{2^4}  \;+\;
      \cdots
      \qquad
      \qquad
      r=\frac{1}{2}
  \]
  \pause
  a série geométrica converge e
  \(\quad \frac{1}{n!} < \frac{1}{2^{n-1}}\quad\)
  para $n>2$

  \vfill
  então a série do inverso do fatorial converge

\end{frame}

%------------------------------------------------------------------------------%
\begin{frame}{Teste da Comparação no Limite}
  Sejam as séries
  \(\; \sum {\color{c1} a_n} \;\) e
  \(\; \sum {\color{c3} b_n} \;\)
  com $\;{\color{c1} a_n} > 0\;$ e $\;{\color{c3} b_n}>0\;$ para todo $n>N$
  \pause
  \vspace{\baselineskip}
  \begin{enumerate}
    \item
    \emph{Se} \(\;\lim \frac{\;{\color{c1} a_n}\;}{{\color{c3} b_n}} = c > 0 \) \;\;
    \emph{então}\; ambas convergem ou ambas divergem

    \pause
    \vfill
    \item
    \emph{Se} \(\;\sum {\color{c3} b_n}\) converge
    \tabto{36mm} \emph{e} \(\;\lim \frac{\;{\color{c1} a_n}\;}{{\color{c3} b_n}} = 0\)
    \tabto{66mm} \emph{então} \(\;\sum {\color{c1} a_n} \) também converge

    \pause
    \vfill
    \item
    \emph{Se} \(\;\sum {\color{c3} b_n}\) diverge
    \tabto{36mm} \emph{e} \(\;\lim \frac{\;{\color{c1} a_n}\;}{{\color{c3} b_n}} = \infty\)
    \tabto{66mm} \emph{então} \(\;\sum {\color{c1} a_n} \) também diverge

    \vfill
  \end{enumerate}
\end{frame}

%------------------------------------------------------------------------------%
\begin{frame}{Exemplo}
  Verificar a convergência da série
  \tabto{60mm}
  \( \sum_{n=1}^{\infty} a_n = \sum_{n=1}^{\infty} \frac{2n+1}{\;n^2+2n+1\;} \)

  \pause
  \vspace{1.7\baselineskip}
  comparamos com
  \tabto{60mm}
  \( \sum_{n=1}^{\infty} b_n = \sum_{n=1}^{\infty} \frac{1}{\;n\;}\qquad \) (Série Harmônica)

  \pause
  \vspace{0.3\baselineskip}
  \[
    \lim_{n\to\infty} \frac{a_n}{b_n}             \pause \;=\;
    \lim_{n\to\infty} \frac{(2n+1)n}{\;n^2+2n+1\;} \pause \;=\;
    \lim_{n\to\infty} \frac{4n+1}{\;2n+2\;}       \pause \;=\;
    \lim_{n\to\infty} \frac{\,4\,}{2}             \;=\; 2
  \]

  \pause
  \vspace{0.5\baselineskip}
  Como a soma de \(b_n\) diverge então a de \(a_n\) também diverge
\end{frame}

%------------------------------------------------------------------------------%
\begin{frame}{\mytitle}{\mysubtitle}
  \tableofcontents
\end{frame}

%------------------------------------------------------------------------------%
\end{document}
%------------------------------------------------------------------------------%
